% cat-third-lift.tex — the third lift: code as a vertex (path::T), and its α-image (abstraction::T).
% Rendered with the tikz MCP server (compile_tikz), which wraps this body in a
% standalone document (default preamble: tikz, amsmath, amssymb + common TikZ
% libraries) and writes docs/website/skills/mtron/assets/cat-third-lift.png.
%
% INK: white on a transparent background — the site's theme is dark (--maincolor
% #000000, panels #191C24). Edge labels are deliberately UNFILLED (no
% `fill=black` patch), so no text box shows against the page; labels are offset
% clear of the strokes instead.
%
% THE FIGURE, IN ONE PICTURE. The third lift: a chain of value code is a path::T
% vertex whose out-edges are rewrite rules (>>r1, >>r2 — rule sites); its
% pointwise α-image is a SINGLE abstraction::T node — the α-word, a reference
% surface with NO edges of its own. α lifts up (a function: one path, one
% abstraction); γ replays a route chosen on the abstraction back at the anchor.
%
%       abstraction::T [A]                                  type code
%       (the α-word —                                        one α-word,
%        the surface a plan reads)                           no edges
%                 ^   :
%            α    |   : γ   replay the route
%           lift  |   :     at the anchor
%       path::T   [p] -- >>r1 --> [p'] -- >>r2 --> [p*]     value code
%       (the member in hand)                                 edges: rewrite rules
%                                                            (rule sites)
%
%                 fiber(A) = { p, q, r, ... }
%
%   * the abstraction is ONE node, not a chain: it pulls nothing up from the
%     concrete chain — it is the α-word (the morphism list) the plan's selector
%     reads to compute cut/span/route. The rewrite edges live on path::T alone.
%   * α (SOLID) lifts UP and is a function: a concrete chain has exactly one
%     abstraction. γ (DASHED) replays DOWN and is a policy: a chosen route is
%     replayed at the anchor, so there is no unique value to draw — the dash is
%     that, not decoration.
%   * Every claim is structural (no counts, no status), so the picture does not
%     age as the third lift lands — transient status stays in the prose.

\begin{tikzpicture}[
  c/.style={draw, rounded corners=2pt, thick, minimum width=10mm, minimum height=6.5mm,
            font=\small, inner sep=1.2mm},
  e/.style={-{Stealth[length=2.2mm]}, thick},
  ed/.style={-{Stealth[length=2.2mm]}, thick, dashed},
  el/.style={font=\scriptsize, inner sep=1pt},
  tl/.style={font=\tiny, inner sep=1pt},
  lbl/.style={font=\small, inner sep=1pt},
  sub/.style={font=\tiny, text=white!60!black, align=left, inner sep=1pt},
  note/.style={font=\scriptsize, inner sep=1pt},
  color=white
]

% ---- the abstraction: ONE node, the α-word (no edges of its own) ----
\node[c] (A) at (0.0,2.6) {$[A]$};

% ---- the path: value code, a chain of rewrite sites ----
\node[c] (p)  at (0.0,0) {$[p]$};
\node[c] (pp) at (3.4,0) {$[p']$};
\node[c] (ps) at (6.8,0) {$[p^{*}]$};
\draw[e] (p)  -- node[el,above=1pt]{$>>\!r_1$} (pp);
\draw[e] (pp) -- node[el,above=1pt]{$>>\!r_2$} (ps);

% ---- the lift: α up (a function — one path, one abstraction) ----
\draw[e] (p.north) -- node[el,left=2pt,pos=0.55]{$\alpha$}
                       node[el,right=2pt,pos=0.55]{lift} (A.south);

% ---- the replay: γ down (a chosen route, replayed at the anchor) ----
\draw[ed] (A.south) -- node[el,right=2pt,pos=0.5]{$\gamma$} (ps.north);
\node[tl, anchor=south, align=center] at (3.3,1.75) {replay the route\\ at the anchor};

% ---- left labels (RIGHT-aligned: the column hugs the boxes) ----
\node[lbl, anchor=east] at (-0.9,2.6)  {\texttt{abstraction::T}};
\node[sub, anchor=east, align=right] at (-0.9,1.95) {(the $\alpha$-word —\\ the surface a plan reads)};
\node[lbl, anchor=east] at (-0.9,0)    {\texttt{path::T}};
\node[sub, anchor=east, align=right] at (-0.9,-0.65){(the member in hand)};

% ---- right labels (LEFT-aligned, off the boxes) ----
\node[lbl, anchor=west] at (7.4,2.6)   {type code};
\node[tl, anchor=west, align=left] at (7.4,1.95) {one $\alpha$-word,\\ no edges};
\node[lbl, anchor=west] at (7.4,0)     {value code};
\node[tl, anchor=west, align=left] at (7.4,-0.65) {edges: rewrite rules\\ (rule sites)};

% ---- the fact under the square ----
\node[note] at (3.2,-1.15) {$\mathrm{fiber}(A) = \{\,p,\,q,\,r,\,\dots\,\}$};
\end{tikzpicture}
